\def\stat{giy}





\def\tit{ПРОБЛЕМЫ ВЗАИМОДЕЙСТВИЯ ВРЕДОНОСНОГО КОДА И~ПРОГРАММ 


ЗАЩИТЫ В~АРХИТЕКТУРЕ СОВРЕМЕННЫХ ОПЕРАЦИОННЫХ  СИСТЕМ$^*$}





\def\titkol{Проблемы взаимодействия вредоносного кода и~программ 


защиты в~архитектуре ОС} %современных операционных систем}





\def\autkol{Р.\,Р.~Гилязов, А.\,А.~Грушо}








\def\aut{Р.\,Р.~Гилязов$^1$, А.\,А.~Грушо$^2$}





\titel{\tit}{\aut}{\autkol}{\titkol}





{\renewcommand{\thefootnote}{\fnsymbol{footnote}}


\footnotetext[1]


{Работа выполнена при частичной финансовой поддержке РФФИ (проект 15-07-02053).}}





\renewcommand{\thefootnote}{\arabic{footnote}}


\footnotetext[1]{Московский государственный университет им.\ М.\,В. Ломоносова, irusrubin@gmail.com}


\footnotetext[2]{Институт проблем информатики Федерального исследовательского


центра <<Информатика и~управление>> Российской академии наук, grusho@yandex.ru}





\label{st\stat}





                  \thispagestyle{headings}





  \Abst{Рассмотрены проблемы взаимодействия вредоносного кода и~защитного 


программного обеспечения (ЗПО) в~средах современных операционных систем (ОС). В~част\-ности, 


рассмотрен ряд аспектов, связанных с~программными модулями, предоставляющими 


нарушителю возможности для устойчивого и~неопределяемого присутствия в~компьютерных 


системах. Сформирован ряд утверждений, говорящих о взаимотношении технологий, 


используемых в~ЗПО, и~обеспечении <<невидимости>> 


исполняемого кода. На практике для модельного руткита (РТ) продемонстрирована возможность 


неопределяемого присутствия для современного ЗПО. На 


основе анализа механизма системных вызовов и~драйверной подсистемы Windows NT 


выработаны необходимые требования по разработке ЗПО. 


Построена модель случайного ограничения возможностей вредоносного программного 


обеспечения (ВПО) для практической реализации ЗПО без нарушения алгоритмов 


распределения процессорного времени.}


   


\KW{информационная безопасность; защитное программное обеспечение; вредоносное 


программное обеспечение; руткит; антивирус; технологии сокрытия исполняемого кода}





\DOI{10.14357\!/\!08696527150306}





                     %\vspace*{-12pt}











\vskip 12pt plus 9pt minus 6pt








      \thispagestyle{headings}





\section{Введение}





  Компьютерные системы, работающие под управлением ОС, 


предоставляют ресурсы для выполнения задач по обработке информации, 


являющейся ценной. Злоумышленник, используя различные методы, имеет 


возможность реализовать специальное программное обеспечение для захвата 


вычислительных ресурсов и~получения несанкционированного доступа к 


ценной информации. Отметим, что нарушитель заинтересован в~постоянном, 


устойчивом и~неопределяемом присутствии своих компонентов. Существует 


класс технологий, позволяющих реализовать данное 


ВПО и~сделать его <<невидимым>> для защитных средств.


  


  Статья посвящена проблеме противодействия защитного программного 


обеспечения ВПО. Для выявления вредоносного кода ЗПО должно 


решать задачи мониторинга, позволяющие выявить ВПО, т.\,е.\ ЗПО должно 


<<увидеть>> ВПО или его следы. Решение этих задач основано на наличии 


двух базовых компонентов в~ЗПО. 


  


  \textbf{Компонент А.} Защитное программное обеспечение должно распознавать следы 


  ВПО, т.\,е.\ должен 


быть язык, на котором выражены те свойства ВПО или следов его 


функционирования, которые система мониторинга способна прочитать, чтобы 


получить информацию о~вредоносном коде. Например, сигнатуры вирусов или 


вредоносного кода являются словами такого языка и~сенсоры ЗПО способны 


идентифицировать эти сигнатуры в~рассматриваемом объекте так, что 


подсистема анализа ЗПО по этой информации способна вывести заключение 


о~наличии вредоносного кода. Если система мониторинга не способна прочитать 


информацию о~том, что данный код является вредоносным, то ЗПО не может 


сделать вывод о~наличии вредоносного кода. В~этом случае обфускация ВПО 


позволяет сделать его <<невидимым>> при статическом анализе. 


  


  \textbf{Компонент В.} Если есть язык, описывающий свойства ВПО и~следов 


его функционирования, то для выявления вредоносного кода должно 


существовать взаимодействие сенсоров ЗПО со статическими или 


динамическими фрагментами ВПО. В~самом деле, любые сенсоры 


регистрируют требуемые параметры по результатам взаимодействия с~объектом 


наблюдения. 


  


  В определенных условиях ВПО может реализовать действия, исключающие 


взаимодействие с~ЗПО. В~этом случае можно также говорить 


о~<<невидимости>> ВПО. Перечислим действия ВПО, позволяющие исключить 


его взаимодействие с~ЗПО. 


  \begin{enumerate}[1.]


\item Использование ограниченности области мониторинга ЗПО. В~этом 


случае ВПО переносит себя в~домен, не доступный для ЗПО. 


\item Вредоносное программное обеспечение использует ограниченность 


времени для обнаружения себя 


в~некоторой среде. Сюда можно отнести возможности ВПО по усложнению 


задачи анализа взаимодействия (например, добавление <<шума>> в~среду 


функционирования ЗПО). 


\item Использование возможности приоритета работы ВПО для действий по 


своей маскировке перед началом функционирования ЗПО путем 


модификации данных ОС, используемых для выявления ВПО. 


\item Использование информации о~последовательности действий ЗПО для 


обхода его защитных функций. 


\item Модификация информации о взаимодействии сенсоров ЗПО с~ВПО. 


Например, перехват IRP-па\-ке\-тов  (I/O request packets) и~их фильтрация.


  \end{enumerate}


  


  В данной статье рассматриваются технологии обеспечения <<невидимости>> 


исполняемого кода в~средах различных ОС. Реализация 


подобных технологий в~программном коде имеет устоявшееся название~--- 


руткит (см., например,~[1]). Актуальными являются вопросы взаимоотношения 


ЗПО и~РТ, связанные с~возможностью его обнаружения. 





\section{Взаимоотношение технологий, используемых в~защитном программном


обеспечении и~руткитах}


  


  Условно РТ можно классифицировать следующим образом~[2]:


  \begin{itemize}


\item стеганографические РТ;


\item РТ, функционирующие в~ОС;


\item РТ, функционирующие вне ОС.


  \end{itemize}


  


  Стеганографические РТ в~данной работе не рассматриваются.


  


  На основе анализа программно-аппаратной архитектуры современных 


ОС в~работе делается следующий вывод. В~случае если ЗПО 


получило процессорные ресурсы раньше РТ и~находится на уровне ниже или 


таком же, как и~РТ, то ЗПО может детектировать~РТ.


  


  Если РТ загрузился раньше ЗПО, он может использовать перечисленные 


выше действия и~либо блокировать загрузку ЗПО, либо создать эмуляцию 


отсутствия РТ относительно ЗПО. Задача идентификации РТ ЗПО идентична 


задаче идентификации ЗПО РТ и~ЗПО ВПО в~целом. Кроме того, на практике 


задача идентификации ЗПО проще ввиду его общедоступности и~известным 


алгоритмам поведения.


  


  На примере ОС двухуровневой архитектуры (Windows NT, Linux, Mac~OS 


и~т.\,п.) рассмотрены модификации источников данных ОС с~целью 


воспрепятствования детектированию РТ ЗПО, функционирующих в~кольце~0. 


Отметим, что создание полной эмуляции в~данной архитектуре ОС в~кольце~0 


для ЗПО невозможно (как минимум для аппаратной архитектуры x86, x86-64) 


по причине того, что в~случае получения ЗПО процессорного времени ЗПО 


получит доступ к процессору и~сможет проанализировать оперативную память, 


что не может быть эмулировано РТ ввиду невозможности виртуализации 


таблицы векторов прерываний (Interrupt Descriptor Table~--- IDT).


  


  Таким образом, для ОС двухуровневой архитектуры можно отметить, что 


ЗПО, работающее в~кольце~0, при условии, что не нарушена его целостность 


и~оно первым получает процессорное время, может обнаружить РТ, в~случае 


если РТ в~момент получения процессорного времени ЗПО присутствует 


в~системе.


  


  Обоснование этого утверждения является крайне важным для подобных 


архитектур, ввиду того что оно опровергает распространенную гипотезу~[3--5], 


утверждающую, что возможно существование РТ, не обнаружимых для ЗПО, 


в~случае более ранней загрузки РТ.


  


  Кроме того, если не используется обфускация, то РТ, хранящийся и/или 


хранящий информацию на жестком диске, детектируется путем сканирования 


жесткого диска средствами ОС и~низкоуровневыми командами с~последующим 


сравнением результатов двух способов. 


  


  В настоящее время используется РТ вне ОС, например возможно 


существование РТ, работающего на уровне гипервизора, работающего под ОС, 


который не может быть обнаружен ЗПО, работающим в~ОС.


  


  Отметим, что необходимы тематические исследования оборудования, 


в~частности процессора, с~целью обнаружения следов функционирования 


программного кода. На текущий момент известен более низкий уровень ВПО 


для процессора по сравнению с~уровнем гипервизора~--- SMM (system management mode)~[6, 7]. Также 


известно, что микрокод процессора можно обновлять~[8] с~целью выявления 


возможных мест хранения РТ (на текущий момент известны техники, 


эксплуатирующие видеоадаптеры, контроллеры жестких дисков~[9], что 


выходит за рамки данной работы, но полученные результаты будут верны и~для 


результатов данных тематических исследований).





\vspace*{-4pt}


  


\section{Построение модельного руткита}





\vspace*{-2pt}





  На примере Windows NT рассмотрим, насколько эффективно ЗПО для данной 


ОС детектирует РТ в~кольце~0. С~этой целью построим модельный РТ, 


использующий указанные выше возможности преодоления ЗПО. 


  


  В Windows NT взаимодействие уровня пользователя и~уровня ядра 


происходит при помощи диспетчера системных вызовов. Для вызова функций 


ядра используется специальная инструкция процессора {\sf sysenter} либо 


прерывание {\sf int~2E}, после чего диспетчер системных вызовов находит 


необходимую функцию по номеру в~одном из двух дескрипторов системных 


сервисов:\linebreak {\sf KeServiceDescriptorTable} и~{\sf KeServiceDescriptorTableShadow}. Каждый 


из них содержит~4~таблицы, из которых в~{\sf KeServiceDescriptorTable} 


используется первая из них, а~в~{\sf KeServiceDescriptorTableShadow} используются 


первая и~вторая. Формат структуры, описывающий каждую таблицу, представлен в листинге~1. 


  


  \renewcommand{\tablename}{\protect\bf Листинг}


  \begin{table}[b]\small %tabl1


  \vspace*{-6pt}


  \begin{center}


  


  


  %\vspace*{2ex}


  


  \begin{tabular}{|l|}


  \hline


  \\[-7pt]


{\sf typedef struct SERVICE\_DESCRIPTOR\_ENTRY}\\


\{\\


{\sf PVOID *Base};\\


{\sf PULONG Count};\\


{\sf ULONG Limit};\\


{\sf PUCHAR Number};\\


\} {\sf SERVICE\_DESCRIPTOR\_ENTRY},\\


{\sf *PSERVICE\_DESCRIPTOR\_ENTRY}\\


\hline


\end{tabular}





\vspace*{-3pt}





\parbox{206pt}{\Caption{Описание структуры\protect\linebreak {\sf KeServiceDescriptorTable}}


  }


\end{center}


\end{table}





  


  Поле {\sf Base} содержит указатель на таблицу адресов функций, которые 


вызываются при обработке системного вызова.


  


  При создании потока поле {\sf ServiceTable} структуры {\sf ETHREAD.Tcb} 


заполняется адресом {\sf KeServiceDescriptorTable} (см.\ листинг~2).


  \begin{table}\small %tabl2


  \begin{center}


    


  \begin{tabular}{|l|}


  \hline


  \\[-9pt]


{\sf nt!KeInitThread+0x65}\\


{\sf 82d3e103 c786bc000000007bb882 mov dword ptr [esi+0BCh]}, \\


    \hspace*{5mm}{\sf offset nt!KeServiceDescriptorTable (82b87b00)}\\


\hline


\end{tabular}





\vspace*{-3pt}








\parbox{274pt}{\Caption{Заполнение поля {\sf ServiceTable} в~функции 


инициализации потока ОС}


}





\end{center}


\vspace*{-3pt}


\end{table}





  \begin{table}[h]\small %tabl3


  \begin{center}





%  \vspace*{2ex}


  


  \begin{tabular}{|l|}


  \hline


  \\[-9pt]


{\sf nt!PsConvertToGuiThread+0x4e}\\


{\sf 82c96165\ c786bc000000407bb882\ mov dword ptr [esi+0BCh]},\\


\hspace*{5mm}{\sf offset\ nt!KeServiceDescriptorTable (82b87b40)}\\


\hline


\end{tabular}





\vspace*{-3pt}





\parbox{276pt}{\Caption{Модификация поля {\sf ServiceTable} потока в~функции {\sf PsConvertToGuiThread}}


}


\end{center}


\vspace*{-9pt}


\end{table}


\begin{table}[b]\small %tabl4


\vspace*{-6pt}


\begin{center}





%\vspace*{2ex}





\tabcolsep=9pt


\begin{tabular}{|l|}


\hline


\\[-9pt]


{\sf dps nt!KeServiceDescriptorTable}\\


{\sf 82bb1b00  82ac643c nt!KiServiceTable}\\


{\sf 82bb1b04  00000000}\\


{\sf 82bb1b08  00000191}\\


{\sf 82bb1b0c  82ac6a84 nt!KiArgumentTable}\\


{\sf dps 82ac643c L 10a}\\


{\sf 82ac643c  82cc1fcf nt!NtAcceptConnectPort}\\


{\sf 82ac6440  82b09855 nt!NtAccessChec}\\


$\ldots$\\


{\sf 82ac6848  82d4cc53 nt!NtQuerySystemEnvironmentValue}\\


{\sf 82ac684c  82d4d247 nt!NtQuerySystemEnvironmentValueEx}\\


{\sf 82ac6850  9dd0c2b0 drv!new\_ZwQuerySystemInformation}\\


{\sf 82ac6854  82cb2174 nt!NtQuerySystemInformationEx}\\


{\sf 82ac6858  82cbfe9e nt!NtQuerySystemTime}\\


{\sf 82ac685c  82d5354e nt!NtQueryTimer}\\


\hline


\end{tabular}








\vspace*{-3pt}





\parbox{274pt}{\Caption{Перечень функций для обработки системных вызовов расположенных 


в~{\sf KeServiceDescriptorTable.Base}}


}


\end{center}


\end{table}


  


  


  В случае если поток использует графическую подсистему, то при первом 


вызове системного вызова подсистемы Win32 происходит вызов функции 


\mbox{{\sf PsConvertToGuiThread}}, которая заменяет значение поля {\sf ServiceTable} на 


\mbox{{\sf KeServiceDescriptorTableShadow}} (см.\ листинг~3).





  


  Реализуем модельный РТ, который модифицирует в~таблице диспетчеризации 


системных служб указатель на функцию {\sf ZwQuerySystemInformation} с~целью 


сокрытия всех процессов в~ОС.


  


  В качестве ЗПО выберем TrendMicro Rootkit Buster~v5.0.0 build~1180, 


последней доступной на момент написания данной работы версии. 


  


  В случае непосредственной модификации соответствующего указателя на 


{\sf ZwQuerySystemInformation} в~{\sf KeServiceDescriptorTable.Base} (см.\ листинг~4)


ЗПО сообщает о~22~угрозах (рис.~1).











  


  \begin{figure} %fig1


      \vspace*{1pt}


 \begin{center}


 \mbox{%


 \epsfxsize=125mm


 \epsfbox{gil-1.eps}


 }


\end{center}


 \vspace*{-11pt}


\Caption{Реакция ЗПО на сокрытие всех процессов в~системе}


\end{figure}


  \begin{table}[b]\small %tabl5


\begin{center}





%\vspace*{2ex}





\tabcolsep=9pt


\begin{tabular}{|l|}


\hline


\\[-9pt]


{\sf NTSTATUS} \\


{\sf PsSetCreateThreadNotifyRoutine}(\\


\hspace*{5mm}{\sf IN PCREATE\_THREAD\_NOTIFY\_ROUTINE NotifyRoutine}\\


);\\


\hline


\end{tabular}





\vspace*{-3pt}





\parbox{286pt}{\Caption{Прототип нотификатора, предоставляемого ОС 


для события создания и~удаления  потока}


}


\end{center}


\vspace*{3pt}


%\end{table}


%\begin{table}\small %tabl6


  \begin{center}





%\vspace*{2ex}





\tabcolsep=26pt


\begin{tabular}{|l|}


  \hline 


  \\[-9pt]


*({\sf PVOID}*)(({\sf PUCHAR}){\sf Ethread}\;+\;{\sf ServieTableOffset})\;=\;\\


{\sf NewKeServiceDescriptorTable};\\


\hline


\end{tabular}





\vspace*{-3pt}





  \parbox{286pt}{\Caption{Пример кода для модификация поля 


  {\sf ServiceTable} структуры {\sf KTHREAD} потока}


}


\end{center}


\end{table}





  Модифицируем технику сокрытия. С~учетом изложенного выше подменим 


указатель значения указателя {\sf ServiceTable} структуры {\sf ETHREAD.Tcb} каждого 


потока, за исключением потоков ЗПО. Для этого зарегистрируем нотификатор на 


создание потока (листинг~5).





  


  Таким образом, при создании каждого потока получаем адрес {\sf ETHREAD} 


и~заменяем поле {\sf ServiceTable} (листинг~6).





\begin{table}\small %tabl7


\begin{center}





%\vspace*{2ex}





\tabcolsep=18pt


\begin{tabular}{|lll|}


\hline


\\[-9pt]


\multicolumn{3}{|l|}{{\sf nt!KiBBTUnexpectedRange:}}\\


{\sf 82a85572 83f910}&{\sf  cmp }    &{\sf ecx,10h}\\


{\sf 82a85575 753d}   &   {\sf jne}&\\


\multicolumn{3}{|l|}{{\sf nt!KiBBTUnexpectedRange+0x42 (82a855b4)}}\\


{\sf 82a85577 52} &       {\sf push}    &   {\sf edx}\\


{\sf 82a85578 53} &     {\sf   push}    &   {\sf ebx}\\


{\sf 82a85579\ e80262da1c}    &{\sf call}&\\


\multicolumn{3}{|l|}{{\sf drv!NewPsConvertToGuiThread (9f82b780)}}\\


{\sf 82a8557e 0bc0}   & {\sf   or}  &   {\sf eax,eax}\\


{\sf 82a85580 58}     &   {\sf pop} & {\sf eax}\\


{\sf 82a85581 5a} &   {\sf pop}&  {\sf       edx}\\


\hline


\end{tabular}





\vspace*{-3pt}





\parbox{256pt}{\Caption{Модификация 


{\sf PsConvertToGuiThread} в~функции {\sf KiBBTUnexpectedRange}}





}


\end{center}


%\end{table}


%\begin{figure} %fig2


%\renewcommand{\figurename}{\protect\bf Рис.}


\setcounter{table}{1}


\renewcommand{\tablename}{\protect\bf Рис.}


\vspace*{1pt}


 \begin{center}


 \mbox{%


 \epsfxsize=112.667mm


 \epsfbox{gil-2.eps}


 }


\end{center}


 \vspace*{-11pt}


\Caption{Модифицированная техника сокрытия}


\end{table}





\renewcommand{\figurename}{\protect\bf Рис.}


\renewcommand{\tablename}{\protect\bf Листинг}





\addtocounter{figure}{1}


\setcounter{table}{7}


  


  


  Для Win32-потоков необходимо найти и~подменить {\sf PsConvertToGuiThread} (см.\ листинг~7).








  В данной функции необходимо выставлять значение {\sf ServiceTable} потока на 


{\sf KeServiceDescriptorTableShadow} (рис.~2).











  На примере диспетчера задач посмотрим, как изменилось значение 


\mbox{{\sf ServiceTable}} в~его потоках (см.\ листинг~8).


  \begin{table}\small  %tabl8


  \begin{center}





%\vspace*{2ex}


  


  \begin{tabular}{|l|}


  \hline


  \\[-9pt]


{\sf !process 0 2 taskmgr.exe}\\


{\sf PROCESS 86668730  SessionId: 1  Cid: 0fb8    Peb:}\\


{\sf 7ffd7000  ParentCid: 01e0}\\


    \hspace*{5mm}{\sf DirBase: 3f021440  ObjectTable: 949e1750}\\


{\sf HandleCount: 112.}\\


    \hspace*{5mm}{\sf Image: taskmgr.exe x}\\


\\


        \hspace*{10mm}{\sf THREAD 86fa1d48  Cid 0fb8.0fbc  Teb: 7ffdf000}\\


{\sf Win32Thread: fe5a6550 WAIT: (WrUserRequest) UserMode Non-Alertable}\\


    \hspace*{15mm}{\sf 85760030  SynchronizationEvent}\\


    \hspace*{10mm}{\sf THREAD 86d4bd48  Cid 0fb8.0fd8  Teb: 7ffd9000}\\


{\sf Win32Thread: 00000000 WAIT: (WrQueue) UserMode Alertable}\\


    \hspace*{15mm}{\sf 86792200  QueueObject}\\


              \\


{\sf dps 86fa1d48+0xbc}\\


{\sf 86fa1e04  86dc65a0}\\


\\


{\sf dps 86dc65a0}\\


{\sf 86dc65a0  867419b8}\\


{\sf 86dc65a4  00000000}\\


{\sf 86dc65a8  00000191}\\


{\sf 86dc65ac  82ac6a84 nt!KiArgumentTable}\\


{\sf 86dc65b0  92e6c000 win32k!W32pServiceTable}\\


{\sf 86dc65b4  00000000}\\


{\sf 86dc65b8  00000339}\\


{\sf 86dc65bc  92e6d02c win32k!W32pArgumentTable}\\


{\sf dps 867419b8 L 10a}\\


{\sf 867419b8  82cc1fcf nt!NtAcceptConnectPort}\\


$\ldots$\\


{\sf 86741dcc  9f82a2b0 drv!new\_ZwQuerySystemInformation} \\


$\ldots$\\


{\sf 86741ddc  82c35874  nt!NtQueryTimerResolution}\\


\hline


\end{tabular}





\vspace*{-3pt}





\parbox{330pt}{\Caption{Модификация поля {\sf ServiceTable} структуры {\sf KTHREAD} потоков 


процесса {\sf taskmgr.exe} с~целью сокрытия процессов}


}


\end{center}


\vspace*{-12pt}


\end{table}





  Таким образом, модификация успешно произошла. 


  


  При применении данной техники реакция ЗПО отсутствует (рис.~3).








  Построим ЗПО, детектирующее РТ кольца~0. Рассмотрим механизм 


сис\-тем\-ных вызовов и~организацию стека драйверов Windows NT для 


определения перечня параметров~\cite{10-g}, которые необходимо проверять 


ЗПО для обнаружения~РТ.





\begin{figure} %fig3


\vspace*{1pt}


 \begin{center}


 \mbox{%


 \epsfxsize=125mm


 \epsfbox{gil-3.eps}


 }


\end{center}


 \vspace*{-11pt}


\Caption{Отсутствие реакции ЗПО на сокрытие всех процессов в~ОС}


\end{figure}











  


  Как было рассмотрено выше, переход на уровень ядра происходит путем 


вызова либо специальной инструкции процессора {\sf sysenter}, либо специально 


зарегистрированного прерывания {\sf int~2Eh}. 


В~современных ОС на базе Windows NT используется только {\sf sysenter}, обработчик 


прерывания {\sf int~2Eh} сохранен для обеспечения обратной совместимости


(см.\ листинг~9).


\begin{table}\small %tabl9


\vspace*{3pt}


\begin{center}





%\vspace*{2ex}





%\tabcolsep=10pt


\begin{tabular}{|lll|}


\hline


&&\\[-9pt]


\multicolumn{3}{|l|}{{\sf ntdll!NtQuerySystemInformation:}}\\


{\sf 77901e20 b89a000000}&{\sf     mov}&    {\sf eax,9Ah}\\


{\sf 77901e25 e803000000}&    {\sf call}&\\


\multicolumn{3}{|l|}{{\sf ntdll!NtQuerySystemInformation+0xd (77901e2d)}}\\


{\sf 77901e2a c21000} &  {\sf  ret} &{\sf 10h}\\


{\sf 77901e2d 8bd4  } &  {\sf  mov}&  {\sf   edx,esp}\\


{\sf 77901e2f 0f34  }     &   {\sf sysenter}&\\


\hline


\end{tabular}





\vspace*{-3pt}





\parbox{224pt}{\Caption{Вызов в~{\sf sysenter} функции\protect\linebreak 


{\sf NtQuerySystemInformation} из библиотеки {\sf ntdll}}


}


\end{center}


\vspace*{-3pt}


\end{table}











     Инструкция {\sf sysenter} использует три {\sf msr}-ре\-гист\-ра процессора~\cite{11-g}:


     \begin{enumerate}[(1)]


\item {\sf IA32\_SYSENTER\_CS}~--- номер {\sf 174h}, устанавливает {\sf cs};


\item {\sf IA32\_SYSENTER\_ESP}~--- номер {\sf 175h}, устанавливает {\sf esp};


\item {\sf IA32\_SYSENTER\_EIP}~--- номер {\sf 176h}, устанавливает {\sf eip}.


\end{enumerate}











\begin{table}\small %tabl10


\begin{minipage}[b]{60mm}%\vspace*{3pt}


\begin{center}


%\vspace*{2ex}


\begin{tabular}{|l|}


\hline


\\[-9pt]


{\sf !idt 2e}\\[-6pt]


\\


{\sf Dumping IDT:\ 80427400}\\[-6pt]


\\


{\sf dd0c68c00000002e:\ 8133c965}\\


{\sf nt!KiSystemService} \\


\hline


\end{tabular}


%\parbox{220pt}


\Caption{Содержание таблицы {\sf IDT} по номеру {\sf 0x2Eh}}


%}


\end{center}


%\vspace*{-3pt}


%\end{table}


\end{minipage}


\hfill


%\begin{table}[h]\small %tabl11


\begin{minipage}[b]{60mm}


  \begin{center}


 %  \vspace*{2ex}


 % \tabcolsep=42pt


  \begin{tabular}{|l|}


  \hline


  \\[-9pt]


{\sf rdmsr\ 176}\\


{\sf msr[176]\;=\;00000000`8133ca50}\\


{\sf u\ 8133ca50}\\


{\sf nt!KiFastCallEntry}\\


\hline


\end{tabular}


 %\parbox{220pt}


 \Caption{Содержание {\sf msr}-регистра с~номером~176}


 \end{center}


 \end{minipage}


 \vspace*{-6pt}


\end{table}





%\renewcommand{\figurename}{\protect\bf Рис.}


%\renewcommand{\tablename}{\protect\bf Листинг}


%\setcounter{table}{11}





  Таким образом, необходимо контролировать корректность обработчика 


прерывания по номеру {\sf 0x2Eh} в~таблице {\sf IDT} (см.\ листинг~10), а~также значение 


{\sf msr}-ре\-гист\-ра с~номером~{\sf 176h}, поскольку при системном вызове регистр {\sf eip} 


устанавливает значение данного регистра (см.\ листинг~11).


  





  Получаем, что {\sf nt!KiFastCallEntry} и~{\sf nt!KiSystemService} являются точками 


входа диспетчера системных вызовов. Диспетчер системных вызовов находит 


необходимую для вызова функцию в~специальных таблицах по 


соответствующим номерам. Данный механизм подробно описан выше. Отсюда 


понятно, что необходимо контролировать значение указателей на функции 


{\sf nt!KiServiceTable} для функций ядра и~{\sf win32k!W32pServiceTable} для 


графической подсистемы.


  


  Проанализировав точку входа в~диспетчер системных вызовов, можно 


заметить, что значение дескриптора системного сервиса берется из структуры 


{\sf ETHREAD.Tcb} текущего потока, которое, в~свою очередь, устанавливается при 


создании потока в~функции {\sf KeInitThread} (см.\ листинг~12).





  \begin{table}[h]\small %tabl12


  \vspace*{3pt}


  \begin{center}





%  \vspace*{2ex}


  


  \begin{tabular}{|lll|}


  \hline


  &&\\[-9pt]


\multicolumn{3}{|l|}{{\sf u nt!KiFastCallEntry\;+\;0xA5 L A}}\\


\multicolumn{3}{|l|}{{\sf nt!KiFastCallEntry+0xa5:}}\\


{\sf 8133caf5 fb }&    {\sf    sti}&\\


{\sf 8133caf6 8bf8   }    &   {\sf mov}&{\sf     edi,eax}\\


{\sf 8133caf8 c1ef08} &   {\sf shr}&   {\sf  edi,8}\\


{\sf 8133cafb 83e710} &{\sf    and}&{\sf     edi,10h}\\


{\sf 8133cafe 8bcf    }   &   {\sf mov}&   {\sf  ecx,edi}\\


{\sf 8133cb00 037e3c }&  {\sf  add}& {\sf    edi,dword ptr [esi+3Ch]}\\


{\sf 8133cb03 8bd8   }&   {\sf mov}&    {\sf ebx,eax}\\


{\sf 8133cb05 25ff0f0000} &  {\sf  and}& {\sf    eax,0FFFh}\\


{\sf 8133cb0a 3b4708 }&   {\sf cmp}&   {\sf  eax,dword ptr [edi+8]}\\


{\sf 8133cb0d 0f83f1f9ffff}&{\sf   jae} &{\sf    nt!KiEndUnexpectedRange  


(8133c504)}\\


\multicolumn{3}{|l|}{{\sf nt!KeInitThread+0x5e:}}\\


{\sf 818f3494 c7473cc0b28e81}&{\sf     mov} &{\sf    dword ptr} \\


\multicolumn{3}{|l|}{{\sf [edi+3Ch],offset nt!KeServiceDescriptorTable (818eb2c0)}}\\


\hline


\end{tabular}


  \parbox{320pt}{\Caption{Использование поля {\sf ServiceTable} 


  структуры {\sf KTHREAD} потока при системном  вызове}}


\end{center}


 \vspace*{-6pt}


\end{table}





  Отсюда получаем, что необходимо контролировать также значение указателя 


{\sf ServiceTable} в~структуре {\sf ETHREAD.Tcb} всех потоков в~системе.


  


  Были рассмотрены возможные точки модификации потока исполнения 


в~реализации системного вызова. Рассмотрим аспекты, связанные с~начальной 


инициализацией данных ядра. Возможные следующие механизмы модификации:


  \begin{itemize}


\item сплайсинг~\cite{10-g} кода диспетчера системных вызовов 


и~функ\-ций-обра\-бот\-чиков;\\[-15pt]


\item модификация указателей на {\sf KeServiceDescriptorTable}\linebreak 


и \mbox{\sf KeServiceDescriptorTableShadow} в~ядре.


\end{itemize}





   Отсюда имеем, что необходимо контролировать целостность ядра путем 


сравнения с~внешним эталоном и~проверять корректность указателей в~ядре на 


{\sf KeServiceDescriptorTable} и~{\sf KeServiceDescriptorTableShadow}.


  


  Итак, был рассмотрен поток исполнения системного вызова либо до ядерной 


структуры, либо до драйвера устройства. Существуют техники 


непосредственной модификации ядерных структур (технологии 


DROM~--- data read only memory~\cite{12-g}), однако полное сокрытие в~данном случае недостижимо 


ввиду необходимости работы с~ними непосредственно ОС. 


Рассмотрим более подробно драйверную подсистему Windows NT.


  


  Драйверная подсистема Windows NT имеет стековую архитектуру. Драйверы 


для устройства могут организовываться в~цепочку и~взаимодействовать друг 


с~другом посредством специального протокола IRP под управлением подсистемы 


ввода/вывода Windows~\cite{10-g}. В~каждом драйвере представлен ряд 


функ\-ций-обработчиков IRP-па\-ке\-тов. Таким образом, получается, что 


необходимо контролировать:


  \begin{itemize}


\item легальность устройств-фильт\-ров, присоединенных к целевому 


устройству;


\item корректность указателей на IRP-обра\-бот\-чи\-ки в~дескрипторе 


драйверов;


\item модификации IRP-обработчиков драйверов.


\end{itemize}





  Суммируя изложенное выше, при отсутствии недекларированных возможностей 


  получаем, что для 


детектирования РТ ЗПО необходимо контролировать следующие параметры:


  \begin{itemize}


\item значения указателей на функции в~{\sf KiServiceTable};


\item значения {\sf msr}-регистров, относящихся к инструкции {\sf SYSENTER};


\item корректность значений {\sf IDT};


\item значение указателя {\sf ServiceTable} в~структуре {\sf ETHREAD.Tcb};


\item корректность указателей на {\sf KeServiceDescriptorTable}\linebreak 


и~\mbox{\sf KeServiceDescriptorTableShadow} в~{\sf KeInitThread} и~{\sf PsConvertToGuiThread};


\item модификации кода диспетчера системных вызовов;


\item модификации функций-обработчиков системного вызова;


\item легальность устройств-фильт\-ров, присоединенных к целевому 


устройству;


\item корректность указателей на IRP-обра\-бот\-чи\-ки в~дескрипторе 


драйверов;


\item модификации IRP-обработчиков драйверов.


\end{itemize}








\section{Модель случайного ограничения возможностей вредоносного программного обеспечения}





  Несмотря на то что теоретически контроль основных параметров ядра ОС 


с~помощью ЗПО возможен, на практике могут возникать большие 


вычислительные трудности. В~связи с~этим возможность полного контроля за 


один квант машинного времени практически не достижима. Отметим, что, 


несмотря на то что ВПО, обладающее достаточными привилегиями, имеет 


воз\-мож\-ность воздействовать на механизм распределения процессорного 


времени, с~практической точки зрения это нецелесообразно, ввиду того что 


оказывает значительное влияние на производительность ОС.


  


  Рассмотрим возможность последовательного контроля основных параметров, 


позволяющих идентифицировать наличие ВПО в~ядре ОС. Предположим, что 


таких параметров~$N$, время дискретно и~на проверку одного параметра 


затрачивается одно и~то же время~$\tau$. Тогда на последовательный просмотр 


всех параметров требуется $N\tau$ времени. Пусть ЗПО способно 


проанализировать корректность одного параметра за один квант машинного 


времени. Если фиксировать порядок проверки, то РТ может со временем 


определить этот порядок и~использовать эту информацию для <<обмана>> ЗПО. 


Единственный способ не допустить воз\-мож\-ности такой атаки РТ на ЗПО~--- это 


случайная генерация номеров контролируемого процесса. Однако построение 


вероятностной модели, позволяющей обойти все контролируемые параметры по 


одному разу в~случайном порядке, достаточно сложно. 


  


  Самой простой схемой является независимый и~равновероятный выбор 


контролируемого параметра на каждом шагу контроля. Такой подход увеличит 


время обхода всех контролируемых параметров, но значительно упростит 


реализацию процесса выбора контролируемого параметра. В~качестве 


генератора случайных чисел можно использовать ка\-кой-ли\-бо случайный 


процесс в~компьютере или генератор псевдослучайных чисел, построенный на 


однонаправленной функции. 


{\looseness=1





}


  


  В~этой модели можно использовать результаты классической задачи 


о~размещениях~[13]. Согласно~[13], время~$v_N$ реализации контроля всех 


необходимых параметров имеет следующие математическое ожидание 


и~дисперсию: 


  \begin{equation}


  {\sf E} v_N = N\sum\limits_{i=1}^N \fr{1}{i}\,;\quad {\sf D}v_N= 


N^2\sum\limits_{i=1}^N \fr{1}{i^2}-{\sf E}v_N\,.


  \label{e1-g}


  \end{equation}


  


  При $N\to\infty$ формулы~(\ref{e1-g}) принимают вид:


  \begin{equation}


  {\sf E}v_N= N\ln N +O(N)\,;\quad {\sf D}v_N=\fr{\pi^2}{6}\,N^2 -N\ln N+O(N)\,.


  \label{e2-g}


  \end{equation}


  


  Из приведенных результатов следует, что использование независимого 


равновероятного выбора очередного параметра контроля приводит 


к~незначительному увеличению времени обхода всех контролируемых 


параметров. 


  


  Данный подход можно обобщить, если за один квант выделяемого 


машинного времени система контроля может 


просмотреть~$k$~контролируемых параметров. При независимом и~


равновероятном выборе $k$ параметров контроля время контроля всех параметров 


увеличивается (см.~(\ref{e2-g})), но возможности ВПО по обману ЗПО 


значительно сокращаются.











\section{Заключение}








  В данной работе были проанализированы технологии взаимодействия 


исполняемого вредоносного кода и~ЗПО. Для ОС двухуровневой архитектуры, 


которые наиболее распространены на текущий момент, обоснован результат, 


состоящий в~том, что если РТ целевым образом не воздействует на ЗПО, то ЗПО 


не всегда может его обнаружить. Несмотря на теоретический характер этого 


результата, практически контроль не достигается в~реализациях современного 


ЗПО. Это демонстрируется в~данной работе путем конструирования модельного 


РТ для одной общедоступной реализации ЗПО. Отсюда следует, что существует 


потребность в~разработке практических правил построения ЗПО, 


обеспечивающих практически необходимый контроль над ВПО. 








  


{\small\frenchspacing


{%\baselineskip=10.8pt


\addcontentsline{toc}{section}{Литература}


\begin{thebibliography}{99}








  \bibitem{1-g}


  \Au{Hoglund G., Butler J.} Rootkits: Subverting the windows kernel.~--- 


Addison-Wesley Professional, 2007. 352~p. 


  \bibitem{2-g}


  \Au{Rutkowska J.} Introducing stealth malware taxonomy~/\!\!/ Help Net Security, 


2006. {\sf http://www.net-security.org/dl/articles/malware-taxonomy.pdf}.


\bibitem{3-g}


\Au{Lau H.} Are MBR infections back in fashion?~/\!\!/ Symantec Corporation, 


2011. {\sf 


http://www.symantec.com/connect/blogs/are-mbr-infections-back-fashion}.


  \bibitem{4-g}


  \Au{Matrosov A., Rodionov E.} Modern Bootkit Trends: Bypassing kernel-mode 


signing policy~/\!\!/ ESET, 2011. {\sf 


  http://www.eset.com/us/resources/white-papers/Rodionov-Matrosov.pdf}.


  \bibitem{5-g}


  \Au{Голованов С., Русаков В.} Атаки до загрузки системы~/\!\!/ 


SECURELIST, 2014. {\sf 


  https://securelist.ru/blog/issledovaniya/20151/ataki-do-zagruzki-sistemy}.


  \bibitem{6-g}


  \Au{Collins R.\,R.} Intel's system management mode. 1997. {\sf 


http://www.rcollins.org/ddj/ Jan97/Jan97.html}.


  \bibitem{7-g}


  \Au{Embleton S., Sparks S.} SMM rootkits. 2008. {\sf 


http://www.hakim.ws/BHUSA08/ speakers/Embleton\_Sparks\_SMM\_Rookits/BH\_US\_08\_Embleton\_Sparks\_ SMM\_Rootkits\_WhitePaper.pdf}.


  \bibitem{8-g}


  Intel(R) 64 and IA-32 Architectures Software Developer's Manual. Vol.~3A: 


System Programming Guide. Pt.~1. Ch.~9.11: Microcode update facilities.~--- 


Intel(R), 2015. P.~432--463. 


  \bibitem{9-g}


  \Au{Menn J.} Russian researchers expose breakthrough U.S.\ spying program. 


2015. {\sf 


  http://www.reuters.com/article/2015/02/16/us-usa-cyberspying-idUSKBN0LK1QV20150216}. 


  \bibitem{10-g}


  \Au{Russinovich M., Solomon D.\,A., Ionescu~A.} Windows internals.~--- 6th 


ed.~--- Microsoft Press Publ., 2012. 752~p.


  \bibitem{11-g}


  Intel(R) 64 and IA-32 Architectures Software Developer's Manual. Vol.~3C: 


System Programming Guide. Pt.~3. Ch.~34.1.~--- Intel(R), 2015. P.~209--210.





\pagebreak





  \bibitem{12-g}


  \Au{Blunden B.} The rootkit arsenal: Escape and evasion in the dark corners of 


the system.~--- Jones\,\&\,Bartlett Learning Publ., 2009. 908~p.


  \bibitem{13-g}


  \Au{Колчин В.\,Ф., Севастьянов Б.\,А., Чистяков~В.\,П.} Случайные 


размещения.~--- М.: Наука, 1976. 224~с.


\end{thebibliography}


} }








\vspace*{-6pt}





\hfill{\small\textit{Поступила в~редакцию 19.08.15}}














\vspace*{9pt}





\hrule





\vspace*{2pt}





\hrule





%\vspace*{9pt}














\def\tit{PROBLEMS OF INTERACTION OF~THE~MALICIOUS CODE 


AND~PROTECTION PROGRAMS IN~ARCHITECTURE OF~MODERN 


OPERATING SYSTEMS}





\def\titkol{Problems of interaction of~the~malicious code 


and~protection programs} % in~architecture of~modern  operating systems}








\def\aut{R.\,R.~Giliazov$^1$ and A.\,A.~Grusho$^2$}





\def\autkol{R.\,R.~Giliazov and A.\,A.~Grusho}











\titel{\tit}{\aut}{\autkol}{\titkol}





\vspace*{-9pt}





\noindent


$^1$M.\,V.~Lomonosov Moscow State University,


1-52 Leninskiye Gory, GSP-1,\linebreak


$\hphantom{^1}$Moscow 119991, Russian Federation





\noindent


$^2$Institute of Informatics Problems,  Federal Research Center ``Computer 


Science\linebreak


$\hphantom{^1}$and Control'' of the Russian Academy of Sciences, 44-2 Vavilov Str.,


Moscow\linebreak


$\hphantom{^1}$119333, Russian Federation








\def\leftfootline{\small{\textbf{\thepage}


\hfill Sistemy i Sredstva Informatiki~---


Systems and Means of Informatics 2015 vol~25 no\,3}


}%


 \def\rightfootline{\small{Sistemy i Sredstva Informatiki~---


 Systems and Means of Informatics   2015   vol~25   no\,3


\hfill \textbf{\thepage}}}





  \Abste{The paper considers the interaction between malware and 


  security software environments of modern operating systems. In particular, 


  a~number of aspects which are related to the software module that provides 


  an opportunity for sustainable and undetectable presence of the offender in 


  computer systems is considered. A~number of statements is made about relationships 


  between the technologies used in security software and ensuring ``invisibility'' 


  of the executable malicious code. The possibility of undetectable rootkit presence 


  in modern security software is shown on practice. In addition, the mechanism 


  of system calls and drivers subsystem of Windows NT is analyzed. Furthermore, 


  necessary practical requirements for implementation of security software are 


  developed. The model of random restriction of malicious software for security 


  software is constructed.} 


  


  \KWE{information security; security software; malicious software; rootkit; 


  antivirus; technologies of hiding execution code}





    \DOI{10.14357\!/\!08696527150306}





%\vspace*{-24pt}








\Ack


\noindent


The paper was partly supported by the Russian Foundation for Basic Research  


(project 15-07-02053).








\renewcommand{\bibname}{\protect\rmfamily References}


%\renewcommand{\bibname}{\large\protect\rm References}





{\small\frenchspacing


{%\baselineskip=10.8pt


\addcontentsline{toc}{section}{References}


\begin{thebibliography}{99}


   \bibitem{1-g-1}


  \Aue{Hoglund, G., and J.~Butler}. 2007. \textit{Rootkits: Subverting the windows 


kernel}. Addison-Wesley Professional Publ. 352~p. 


  \bibitem{2-g-1}


  \Aue{Rutkowska, J.} 2006. Introducing stealth malware taxonomy. Help net 


security. Available at: {\sf 


  http://www.net-security.org/dl/articles/malware-taxonomy.pdf} (accessed 


August~17, 2015).


  \bibitem{3-g-1}


  \Aue{Lau, H.} 2011. Are MBR infections back in fashion? 


\textit{Symantec Corporation}. Available at: {\sf 


http://www.symantec.com/connect/blogs/are-mbr-infections-back-fashion} (accessed 


August~17, 2015).


  \bibitem{4-g-1}


  \Aue{Matrosov, A., and~E.~Rodionov}. 2011. Modern Bootkit Trends: Bypassing 


kernel-mode signing policy. \textit{ESET}. Available at: {\sf 


http://www.eset.com/us/resources/white-papers/Rodionov-Matrosov.pdf} (accessed 


August~17, 2015).


  \bibitem{5-g-1}


  \Aue{Golovanov, S., and V.~Rusakov}. 2014. Ataki do zagruzki sistemy 


[Attacks before loading of system]. \textit{SECURELIST}. Available at: {\sf 


https://securelist.ru/blog/ issledovaniya/20151/ataki-do-zagruzki-sistemy/} (accessed 


August~17, 2015).


  \bibitem{6-g-1}


  \Aue{Collins, R.\,R.} 1997. Intel's system management mode. Available at: {\sf 


http:// www.rcollins.org/ddj/Jan97/Jan97.html} (accessed August~17, 2015).


  \bibitem{7-g-1}


  \Aue{Embleton, S., and S.~Sparks}. 2008. SMM rootkits.  Available at: 


{\sf 


http:// www.hakim.ws/BHUSA08/speakers/Embleton\_Sparks\_SMM\_Rookits/BH\_US\_08\_ Embleton\_Sparks\_SMM\_Rootkits\_WhitePaper.pdf} (accessed 


August~17, 2015).


  \bibitem{8-g-1}


  Intel(R). 2015. Intel(R) 64 and IA-32 Architectures Software Developer's Manual. Vol.~3A: 


System Programming Guide. Pt.~1. Ch.~9.11: Microcode update facilities. 432--463. 


  \bibitem{9-g-1}


  \Aue{Menn, J.} 2015. Russian researchers expose breakthrough 


U.S.\ spying program. Available at: {\sf 


http://www.reuters.com/article/2015/02/16/us-usa-cyberspying-


idUSKBN0LK1QV20150216} (accessed August~17, 2015). 


  \bibitem{10-g-1}


  \Aue{Russinovich, M., D.\,A.~Solomon, and A.~Ionescu}. 2012. \textit{Windows 


internals}. 6th ed. Microsoft Press Publ. 752~p.


  \bibitem{11-g-1}


  Intel(R). 2015. Intel(R) 64 and IA-32 Architectures Software Developer's Manual. Vol.~3C: 


System Programming Guide. Pt.~3. Ch.~34.1. 209--210.


 


  \bibitem{13-g-1}


  \Aue{Blunden, B.} 2009. \textit{The rootkit arsenal: Escape and evasion in the 


dark corners of the system}. Jones\,\&\,Bartlett Learning Publ. 908~p.


  \bibitem{14-g-1}


  \Aue{Kolchin, V.\,F., B.\,A.~Sevastyanov, and V.\,P.~Chistyakov}. 1976. 


\textit{Sluchaynye razmeshcheniya} [Random placements].~--- M.: Nauka. 224~p.


  \end{thebibliography}


} }








\vspace*{-1pt}





\hfill{\small\textit{Received August 19, 2015}}  








\Contr





\noindent


\textbf{Giliazov Ruslan R.} (b.\ 1989)~--- junior scientist, Faculty of Computational Mathematics 


and Cybernetics, M.\,V.~Lomonosov Moscow State University, 1-52 Leninskiye Gory, GSP-1, 


Moscow 119991, Russian Federation; irusrubin@gmail.com





\vspace*{5pt}





\noindent


\textbf{Grusho Alexander A.} (b.\ 1946)~--- Doctor of Science in physics and mathematics, 


professor, leading scientist, Institute of Informatics Problems,  Federal Research Center ``Computer 


Science and Control'' of the Russian Academy of Sciences, 44-2 Vavilov Str.,


Moscow 119333, Russian Federation; grusho@yandex.ru








 \label{end\stat}








\renewcommand{\tablename}{\protect\bf Таблица}





\renewcommand{\bibname}{\protect\rm Литература}


